%A LaTeX model for SYSU General Phys. Lab. by Probfia Gao.
%用XeLaTeX编译
\documentclass[11pt,a4paper]{ctexart}

%在下面补全实验名，例如 实验BB3 光电效应实验。
\newcommand{\ExpeName}{实验D1 锁相放大器与弱信号测量}
%学生名字和学号
\newcommand{\Stuname}{高寒}
\newcommand{\StuID}{17353019}

\usepackage{fancyhdr}
\usepackage{amsmath}
\usepackage{graphicx}
\usepackage[hmargin=1.25in,vmargin=1in]{geometry}
\usepackage{pdfpages}

\usepackage[colorlinks,
            linkcolor=red,
		 urlcolor=black]{hyperref}

\usepackage{float}
\usepackage{cleveref}

\crefname{equation}{}{}
\crefname{figure}{图}{图}
\crefname{footnote}{注释}{注释}
\crefname{table}{表}{表}

%\cpic{<尺寸>}{<文件名>}}用于生成居中的图片。
\newcommand{\cpic}[2]{
\begin{center}
\includegraphics[scale=#1]{#2}
\end{center}
}

%\cpicn{<尺寸>}{<文件名>}{<注释>}用于生成居中且带有注释的图片，其label为图片名。
\newcommand{\cpicn}[3]
{
\begin{figure}[H]
\cpic{#1}{#2}
\caption{#3\label{#2}}
\end{figure}
}

\newcommand{\red}{\color{red}}

\newcommand{\beq}{\begin{equation}}
\newcommand{\eeq}{\end{equation}}
\newcommand{\bea}{\begin{equation}\begin{aligned}}
\newcommand{\eea}{\end{aligned}\end{equation}}

%输入单位和数学常数
%下面所有命令需在公式环境下使用
\newcommand{\e}{\mathrm{e}}   %自然常数e = \e
\newcommand{\im}{\mathrm{i}}   %虚数单位i = \im
\newcommand{\meter}{\mathrm{m}}      %单位/前缀 = \单位/前缀英文名
\newcommand{\newton}{\mathrm{N}}  
\newcommand{\joule}{\mathrm{J}}
\newcommand{\second}{\mathrm{s}}
\newcommand{\gram}{\mathrm{g}}
\newcommand{\ampere}{\mathrm{A}}
\newcommand{\kilo}{\mathrm{k}}
\newcommand{\milli}{\mathrm{m}}
\newcommand{\kelvin}{\mathrm{K}}
\newcommand{\mole}{\mathrm{mol}}
\newcommand{\volt}{\mathrm{V}}
\newcommand{\nano}{\mathrm{n}}
\newcommand{\degreeC}{^\circ \mathrm{C}}  %摄氏度符号 = \degreeC

\newcommand{\unit}[1]{\rm \ #1}
\newcommand{\emptyline}{\par \ \\ }

\pagestyle{fancy}

\fancyhead[L]{\footnotesize{中山大学物理与天文学院近代物理实验}}
\fancyhead[R]{\footnotesize{\ExpeName}}
\fancyfoot[C]{\thepage}

\begin{document}
%第一页
\cpic{0.255}{e1}%学生信息和计分表格
\begin{table}[H]
\centering
\begin{tabular}{|p{32mm}|p{32mm}|p{32mm}|p{32mm}|}
\hline
年级、专业： & 17级 物理学 & 组号： & 实验班6 \\ \hline
姓名： & \Stuname & 学号： & \StuID \\ \hline
日期： & \today & 教师签名： &  \\ \hline
\end{tabular}
\end{table}
\begin{center}
{\LARGE\textbf{{\ExpeName}}}
\end{center}
{\large \bf【实验报告注意事项】}
\begin{enumerate}
 \item 实验报告由三部分组成：
 \begin{enumerate}
  \item[1)]预习报告：（提前一周）认真研读\textbf{\uline{实验讲义}}，弄清实验原理；实验所需的仪器设备、用具及其使用（强烈建议到实验室预习），完成讲义中的预习思考题；了解实验需要测量的物理量，并根据要求提前准备实验记录表格（由学生自己在实验前设计好，可以打印）。预习成绩低于50\%者不能做实验{\color{red} (实验D2和D3需要提前一周的周四完成预习报告交任课老师批改，批改通过后，才允许做实验)}。

  \item[2)]实验记录：认真、客观记录实验条件、实验过程中的现象以及数据。实验记录请用珠笔或者钢笔书写并签名（{\color{red}用铅笔记录的被认为无效}）。{\color{red}保持原始记录，包括写错删除部分，如因误记需要修改记录，必须按规范修改。}（不得输入电脑打印，但可扫描手记后打印扫描件）；离开前请实验教师检查记录并签名。
  \item[3)]分析讨论：处理实验原始数据（学习仪器使用类型的实验除外），对数据的可靠性和合理性进行分析；按规范呈现数据和结果（图、表），包括数据、图表按顺序编号及其引用；分析物理现象（含回答实验思考题，写出问题思考过程，必要时按规范引用数据）；最后得出结论。
 \end{enumerate}
 \textbf{实验报告}就是预习报告、实验记录、和数据处理与分析合起来，加上本页封面。
 \item 每次完成实验后的一周内交\textbf{实验报告}。
 \item 除实验记录外，实验报告其他部分建议双面打印。
\end{enumerate}

\ 
\\
\ 

\begin{flushright}                                                           %模板作者
\tiny{
A \LaTeX \ model for Modern Phys. Lab., SPA, SYSU by \em{\href{https://www.weibo.com/3532532974/profile?rightmod=1&wvr=6&mod=personinfo&is_all=1}{Probfia} Gao.}
}
\end{flushright}

\newpage
{\large{\bf【实验安全注意事项】}}{\red
\begin{enumerate}
\item OE1022 锁相放大器输入端不能接入强信号;接入信号需限定在$1\unit V$以下。
\item OE4004 稳流源的电流输出不能直接接入 OE1022 锁相放大器的输入端。
\end{enumerate}
}

\newpage
\newpage%预习报告
\begin{center}
{\LARGE{\textbf{\ExpeName}}}
\end{center}
\textbf{【实验目的】}
\begin{enumerate}
 \item[1.]了解锁相放大器工作原理和特点，掌握锁相放大器基本参数含义及锁相放大器的基本 操作;复习示波器的使用；（实验内容 1）
  \item[2.]掌握用锁相放大器检测出湮没于噪声中的弱信号方法；（实验内容 2）
  \item[3.]通过本实验更加直观的了解方波的高频成分，并掌握用锁相放大器检测微弱信号中的谐波；（实验内容 3）
\item[4.]掌握交流四引线法精密测量小阻抗的原理和方法，理解相位的物理含义及电压与电流之间相位差的测量方法；（实验内容 4，探索）
\item[5.] 学习 VISA 接口协议，通过 LabVIEW 环境用 PC 机控制锁相放大器数据采集；（实验内容 4）
\item[6.]了解 PN 结结电容，探究变容二极管内 PN 结电容与反偏电压的关系；（实验内容 5）
\item[7.] 通过实验了解热噪声的特征，理解热噪声的本质及其统计意义。（实验内容 6，探索）
\end{enumerate}
\emptyline
\\
\textbf{【仪器用具】}
%将讲义中的表格截图保存为t1在该文件夹下后删去下一行之前的%符号，合理调整scale参数。
\begin{table}[H]
\cpic{0.4}{t1}
\caption{弱信号检测实验仪器用具}
\end{table}
%或者自己去 https://www.tablesgenerator.com/ 做一个表。
\emptyline
\\
\textbf{【原理概述】}\par
现代测量中，所有的物理最终都转换为电压或光强进行记录和处理，我们称之为携带被测量物理量信息的电或光信号；然而所有的测量，即使完全由机器自动进执行，结果都不可避免地引入待测量以外的其他信息，如环境干扰、不确定等，在所测量的信号不可避免地携带了噪声\cite{pulch2013sensitivity,error}。测量结果（信号$x(t)$）可以视作被测量信号$s(t)$与噪声$n(t)$的叠加$x(t) = s(t) + n(t) $。
\par
一般来说，信号是运载信息的工具，是信息的载体。任何携带某个现象属性或行为信息的物理量都可以作为信号，如光信号、声信号和电信号等。物理世界中的信号都是有规律地随时间或空间变化的。数学上以函数的形式$x(t)$描述信号。
\par
最简单的信号用单一频率的三角函数的描述
\beq \label{eq1}
s_\omega (t) = a \cos (\omega_0 t + \varphi) 
\eeq
它包含了三个最基本的信息，即振幅$a$、频率$\omega_0$和相位$\varphi$，它对于信号载体对应着不同的物理含义。在实际应用中，人们往往只用到其中的一个或两个基本信息来进行信号传递或处理，而对于科学探索可能会全方位地用到所有的信息。\cref{eq1}是在时域下对信号的描述;在频域下该信号的数学描述为：
\beq
s_e(\omega) = \begin{cases}
a &, \omega = \omega_0 \\
0 &, \omega \not= \omega_0
\end{cases}
\eeq
其形式更简单，但丢失了相位的信息。基本信号的更一般数学描述是复数形式：
\beq
s_\omega (t) = ae^{i (\omega t + \phi)}
\eeq
它包含了式$\cref{eq1}$描述的实况和虚部，具有更丰富的物理内含和更简洁的数学表达。
\par
一般周期性的信号（函数）都可以傅立叶级数展开：
\beq \label{eq4}
s(t) = \sum_n a_n e^{i(n \omega_0 t + \varphi)}
\eeq
因此，只要讲清楚一种基本形式的信号处理过程，就可推广到一般的信号处理过程。
\par
广义地说，不是待测信号本身的、对测量值的贡献都可视为噪声。它可分为来自外界的环境噪声，最典型的是市电噪声（$50\unit{Hz}$）;以及来自被测量对象本身的噪声，如热噪声。它们都混合待测信号的测量值中。
\par
从物理角度来看，任何待物理量都包含无规、随机的变化，这就是噪声。它是对有规律的信号的一种随机干扰，往往被扩展为与待测信号相冲突的无用信号。例如白噪声，是电子器件和电路中常见的一种噪声，电阻的热噪声，PN结的散弹噪声等都是白噪声，它在时域和频域都呈均匀分布。换句话说，{\red 对于一项实验，噪声可来自于实验对象本身，也可来自于测量系统， 包括传感器和测量仪器，更可能来自于实验系统以外的环境}。
\par
对于有规律变化的环境噪声，可用\cref{eq4}式来表述：它的频率可能与待测信号的频率相同或相近，但更可能是其他频率。在时域可以看成频率不同、振幅随机的“信号”叠加\footnote{单个傅立叶模的振幅不应显含时间$t$，讲义拟有误。}：
\beq
n(t) = \sum_\omega b(\omega) e^{i\omega t}
\eeq
其中振幅$b(\omega)$可能是常量（如市频干扰），亦可能是随机量$r(\omega)$；在频域中，随机噪声也可以理解为每一时刻的噪声电压频谱$n(\omega) $都不会是一样的\footnote{拟有误，随机噪声应该理解成一维随机高斯场。}。
\par
以上讲义内容都是鬼扯，要理解噪声必须明白随机过程的概念和两点关联函数，功率谱等关键统计量。一段信号$f(t)$的两点关联函数\cite{prob}定义为
\beq
\xi(t_2,t_1) = \langle f(t_2) f(t_1) \rangle
\eeq
根据噪声随时间均匀同性的特点，单点关联函数，即噪声在某点上的期望$\langle \xi(t)\rangle$必须为0；$\xi(t_2,t_1)$必须只依赖于$|t_2 - t_1|$。记$\xi(|t_2-t_1|) = \xi(t)$，可以求$\xi(t)$的傅立叶变换
\bea
\xi(\omega_1,\omega_2) &= \int dt_1 dt_2 e^{i\omega_1 t_1 + i\omega_2 t_2} \xi(t_2,t_1) \\
&= \int dt dt_2 e^{i\omega_1 (t_2 + t) + i\omega_2 t_2} \xi(t) \\
&= 2\pi \delta(\omega_1 +\omega_2) \int dt e^{i\omega_1 t}   \xi(t)
\eea
其中引入了记号$t = t_1 - t_2$和积分公式
\beq
\int dt e^{i \omega t} = 2\pi \delta(\omega)
\eeq
积分$\int dt e^{i\omega_1 t}   \xi(t)$的结果只依赖于模式$\omega_1$，定义功率谱
\beq \label{powsp}
P(\omega)= \int dt \xi(t) e^{i\omega t}
\eeq
进一步利用$\xi(t) = \xi(-t)$的性质写出
\beq
P(\omega) = \int_{t>0} \xi(t) \cos \omega t
\eeq
可以立刻看出$P(\omega) = P(-\omega)$，于是实际问题中通常只讨论$\omega>0$的部分。\par
现在傅立叶空间中的两点关联函数\cite{cosmo,dtsft}可以写成
\beq
\xi(\omega_2,\omega_1) = 2\pi \delta(\omega_1 + \omega_2) P(\omega_1)
\eeq
变换回时域空间，两点关联函数
\bea
\xi(t_2,t_1) &= \int \frac{d\omega_2 d\omega_1}{(2\pi)^2} 2\pi \delta(\omega_1 + \omega_2) P(\omega_1) e^{i\omega_1 t_1 + i\omega_2 t_2} \\
&= \int \frac{d\omega}{2\pi} P(\omega) e^{i\omega(t_1 - t_2)}
\eea
特别地，信号在某点处的均方
\beq \label{rms}
\langle f(t)^2 \rangle = \xi(t,t) = \int \frac{d\omega}{2\pi} P(\omega)
\eeq
这也是$P(\omega)$被称作功率谱的原因。对应完全的噪声信号，上式就等于方差。
\par
功率谱是信号的基本性质，在宇宙学中有非常重要的作用，例如对于不同的暴涨模型可以通过计算暴涨期间量子涨落带来的功率谱得到暴涨结束时的功率谱，最后通过功率谱计算晚期宇宙结构演化验证暴涨模型的正确与否\cite{cosmo}。
\par
现在我们可以给出白噪声的严格定义：白噪声定义为功率谱为常数的噪声信号。
\par
信噪比定义为信号的有效值（均方根）和噪声的有效值之比
\beq
SNR = \frac{S}{N}
\eeq
得到信噪比的频域表达是有用的。设$f(t)$为一段信号，按傅立叶模展开为
\beq
f(t) = \int \frac{d\omega}{2\pi} f(\omega) e^{i\omega t}
\eeq
其平方积分为
\bea
\int dt [f(t)]^2 &= \int dt \int \frac{d\omega_1}{2\pi} f(\omega_1) e^{i\omega_1 t} \int \frac{d\omega_2}{2\pi} f(\omega_2) e^{i\omega_2 t}  \\
&= \int \frac{d\omega_1 d\omega_2}{(2\pi)^2}  f(\omega_1)f(\omega_2)  2\pi \delta(\omega_1 + \omega_) \\
&= \int \frac{d\omega}{2\pi} f(\omega) f(-\omega) = \int \frac{d\omega}{2\pi} |f(\omega)|^2
\eea
信噪比是科学与工程中所用的一种度量，用于比较所需信号的强度与背景噪声的强度，表 征噪声对信号的覆盖程度。
通过信号处理可提高信噪比，即提高了信号质量。为反映信号质量改善的程度，定义信噪改善比$SNIR$。
\beq
SNIR = \frac{SNR_0}{SNR_i}
\eeq
式中，$SNR_o$是经过处理（系统输出端）的信噪比，$SNR_i$是未经过处理（系统输入端）的信噪 比。$SNIR$越大，表明系统抑制噪声的能力越强。
\par
“微弱信号”不只意味着信号的绝对幅度很小，而是相对于噪声而言幅度很小，特别是那些被噪声淹没的信号，都称作微弱信号。只有在有效地抑制噪声的条件下放大微弱信号的幅度， 才能提取出真实的信号。因此，微弱信号检测所一门专门针对噪声的技术，其主要任务是提高信噪比。\par
提高信噪比的途径有两方面：一是通过物理屏蔽来抑制环境噪声，还可以通过不做实验来避免遭遇噪声，另一方面是通过信号处理滤除噪声。通常物理屏蔽的代价不少，关键是有另一类噪声在理论上不能通过物理屏蔽来抑制或隔断，它就是来自测量对象本身的噪声，尽管此类噪声在幅度上可能远低于环境噪声；因此，对于不能被屏蔽的噪声，信号处理就成为唯一的选择。\par
为了从噪声中提取出有用的信号，就需要研究噪声的来源与性质，分析噪声产生的原因和规律，以及噪声的传播途径，有针对性地采取有效措施抑制噪声，研究被测信号和噪声的统计特性及其差别，以寻找出从背景噪声中检测出有用信号的方法。
\par
小信号是相对微弱信号信噪比更高的信号，可以采用{\bf 滤波}的方法进行信噪比的提升以实现精确测量。这个说白了就是把不需要的傅立叶模式删掉。
\par
设输入信号为$s(t) + n(t)$，信号$s(t)$的傅立叶模单一，为$a \delta(\omega - \omega_0)$；噪声的傅立叶变换由其功率谱决定，傅立叶模式的涨落幅度为$\sim \sqrt{P(\omega)}$，其中$P(\omega)$为前面提到的功率谱。总输入的傅立叶变换为
\beq
f(\omega) = a \delta(\omega - \omega_0) + \sqrt{P(\omega)}
\eeq
其中$\sqrt{P(\omega)}$表示均值为0，方差为$P(\omega)$的一个随机变量，其每个时刻的值是未知的，但功率可以对功率谱积分算出。带通滤波的作用是删去$(\omega_0 - \frac{\Delta \omega}{2} , \omega_0+ \frac{\Delta \omega}{2})$以外的信号，滤波后的总信号为
\bea
f_o (t) &= \int_{\omega - \frac{\Delta \omega}{2}}^{\omega+ \frac{\Delta \omega}{2}} \frac{d\omega}{2\pi} \left(a \delta(\omega - \omega_0) + \sqrt{P(\omega)}\right) e^{i \omega t} \\
&= ae^{i\omega_0 t} + \int_{\omega - \frac{\Delta \omega}{2}}^{\omega+ \frac{\Delta \omega}{2}} \frac{d\omega}{2\pi} \sqrt{P(\omega)}e^{i \omega t}
\eea
滤波前后的信噪比分别为
\beq
SNR_i = \frac{a^2}{\int_{-\infty}^\infty \frac{d\omega}{2\pi} P(\omega)}
\eeq
\beq
SNR_o = \frac{a^2}{\int_{\omega_0 - \frac{\Delta \omega}{2}}^{\omega_0 + \frac{\Delta \omega}{2}} \frac{d\omega}{2\pi} P(\omega)}
\eeq
信噪改善比
\beq
SNIR = \frac{SNR_o}{SNR_i} = \frac{\int_{-\infty}^{\infty} d\omega P(\omega)}{\int_{\omega_0 - \frac{\Delta \omega}{2}}^{\omega_0+ \frac{\Delta \omega}{2}}    d\omega P(\omega)}
\eeq
例如当噪声信号为分布宽度为$\Delta \omega_0$的白噪声时
\beq
SNIR = \frac{\Delta \omega_0}{\Delta \omega} > 1
\eeq
\par
低通滤波是将频率高于$\omega_c$的频率作截断，再将剩余的信号放大$A$倍。\footnote{再次强调，$\sqrt{P(\omega)}$不是确定的傅立叶模大小，而是一个随机变量。}
\beq
f(t) = Aae^{i\omega_0 t} + A \int_{-\infty}^{\omega_c} \sqrt{P(\omega)} e^{i \omega t}
\eeq
当然我们需要$\omega_c > \omega_0$。和之前的推导无异，现在的信噪改善比
\beq
SNIR = \frac{\int_{-\infty}^{\infty} d\omega P(\omega)}{\int_{-\infty}^{\omega_c}    d\omega P(\omega)}
\eeq
由于$P(\omega)$半正定，上式恒大于1。
\par
实际中可以用示波器进行简单滤波，如\cref{filter}。
\cpicn{0.6}{filter}{$Q$值为 20 的滤波器滤波前后的信号:(a)、(b)分别为滤波前后的时域图;(c)、(d)分别为滤波前后的频域图。}
\par
另一种削弱噪声信号的方法为{\bf 零位法}，零位法的本质是比较法测量。一般直接指示型测量仪表的测量方法是将被测信号放大到一定幅度，以驱动表头指针的偏转角度指示被测量的大小;或者经模数转换和数据处理后由液晶显示被测量的数值。而零位法是调整对比量的大小使其尽量接近被测量，由对比量指示被测量的大小，如\cref{d16}所示。图中的零位表指针只用来指示被测量和对比量的差异值，当零位表指示近似为零时，对比量的大小就表征了被测量的大小。对比量的调整可以手动实现，也可闭 环自动调整，如\cref{d16}中虚线所示。用这种方法测量的分辨率取决于对比量的调整和指示的分辨率。用机械测量类比，弹簧秤是一种直接指示仪表，而天平则是一种零位法仪表。
\cpicn{0.6}{d16}{零位法测量原理}
\par
由\cref{d16}可以看出，虽然被测量和对比量在传输过程中分别附加了干扰噪声$n_1(t)$和$n_2(t)$，但是在对比相减的过程中，$n_1(t)$和$n_2(t)$会在一定程度上相互抵消。两路信号传输过程越相似，$n_1(t)$和$n_2(t)$也会越近似，抵消作用越好。因此，与直接指示测量方法相比，零位法测量结果的信噪比更高，测量精度也更高。
\par
设被测信号
\beq
f(t) = s(t) + n_1 (t)
\eeq
对比量信号
\beq
f_r(t) = s_r(t) + n_2 (t)
\eeq
测量其差值得到
\beq
\Delta f (t) = s(t) - s_r(t) + n_1(t) - n_2(t)
\eeq
很容易想到，调节$s_r(t)$使得$\Delta f(t)$的功率谱最小时，$s(t) = s_r(t)$\footnote{讲义上原为$\Delta f(t) = 0$，这是显然错误的。}设两个噪声信号完全相干，那么噪声的功率谱将相消而降为零；若两个噪声完全不相干，其差值的两点关联函数
\beq
\langle (n_1(t_1) - n_2(t_1))(n_1(t_2)-n_2(t_2))\rangle = \langle n_1(t_1) n_1(t_2) \rangle + \langle n_2(t_1) n_2(t_2)\rangle
\eeq
故其功率谱为两个噪声功率谱相加。虽然也可以调节$s_r$使得总功率最小来测量原信号，但如果噪声信号非常强，在操作上很难进行，因此这种方法只适合相干噪声的消除。
\par
本实验重点关注锁相放大法。锁相放大器也称锁定放大器，是通过特定的载波将湮没在噪声中的微弱信号提取出来的技术。锁相放大技术采用了频谱迁移和交流放大技术、用相敏检测器和低通滤波器来实现调制信号的解调。锁相放大器抑制噪声有三个基本出发点:
\begin{itemize}
\item 用调制器将直流或慢变信号的颏谱迁移到调制频率$\omega_0$处,再进行放大,以避开传输过程所引入的以及放大器内部的$\frac{1}{f}$噪声的不利影响。
\item 利用相敏检测器实现被调制信号的解调过程,可以同时利用频率$\omega_0$和相角$\theta$进行检测,噪声与信号同頻又同相的概率很低。
\item 用低通滤波器来抑制宽带噪声。电路实现时，低通滤波器的频带可以做得很窄,而且 其频带宽度不受调制频率的影响。
\end{itemize}
锁相放大技术需要先对（慢变）信号频谱进行迁移，如\cref{d17}，其调制过程是将低 频信号$V$乘以频率为$\omega_0$的正弦载波,从而将其频谱迁移到调制频率$\omega_0$的附近,之后进行选频放大，这样就不会把$\frac{1}{f}$噪声和低频漂移也放大了,如\cref{d17}(a)所 示，图中的虚线表示$\frac{1}{f}$噪声和白噪声的功率谱密度。经交流放大后，再用相敏检测器将其频谱迁移到直流（$\omega = 0$）的两边，用窄带低通滤波器滤除噪声，在这里加一句废话应该看不见，就得到高信噪比的放大信号，如\cref{d17}(b)所示,图中用虚线表示低通滤波（LPF）的频率响应曲线。只要 LPF 的带宽足够窄,就能有效地改善信噪比。
\cpicn{0.6}{d17}{锁相放大器对信号频谱进行迁移的过程}
通过锁相放大技术的过程，相当于一个$Q$值极高的窄带滤波器，可以不受待测频率以外的噪声信号影响，从而使检测系统的性能大为改善。锁相放大器的等效噪声带宽可以达到$0.0001\unit{Hz}$，整体增益可以高达 1010 以上，即把$0.1\unit{nV}$的微弱信号可以放大到$1.0\unit{V}$以上，因而可以测量极其微小的信号。
\emptyline
\\
\textbf{【工作原理】}\par
锁相放大器的基本结构如\cref{d19}所示的虚线框内，其中信号通道、参考通道为锁相放大器的输入通道，相敏检测器（PSD）和低通滤波器（LPF）等。对于三角函数信号，可直接从信号通道输入（虚线框内的）锁相放大器。对于非三角函数信号或慢信号（如直流信号），它在输入锁相放大器前，需要被与参考信号频率相同的信号所调制（如\cref{d19}虚线框外部分所示)，参考信号既可以是外部输入，也可以是内部自带参考信号源提供。为不失一般性，以下数学推导以需要调制的非三角函数信号为出发点。
\cpicn{0.4}{d19}{锁相放大器工作原理（虚线框内）}
\par
设带噪声的待测信号
\beq
f(t)= s(t) + n(t)
\eeq
调制并放大$A_I$倍后的信号为
\beq
u(t) = A_I f(t) \sin \left(\omega_{0} t+\theta\right) = A_I s(t) \sin \left(\omega_{0} t+\theta\right) + A_I n(t) e\sin \left(\omega_{0} t+\theta\right)
\eeq
此时，信号与噪声都被同时放大了。
\par
锁相放大器采用参考信号$u_{r}(t)=\sin \left(\omega_{0} t\right)$解调，因调制信号是用参考信号通过某种物理机制触发或用参考信号本身，因此，参考信号与调制信号频率相同、相位差确定。
\bea
u_{p}(t)&=u_{r}(t) u_{\mathrm{s}}(t)=A_{I}[s(t)+n(t)] \sin \left(\omega_{0} t+\theta\right) \sin \left(\omega_{0} t\right) \\ &=\frac{1}{2} A_{I}[s(t)+n(t)]\left[\cos \theta-\cos \left(2 \omega_{0} t+\theta\right)\right]
\eea
低通滤波后留下了慢变的待测信号和同（低）频噪声：
\beq
u_{ox}(t)=\frac{1}{2} A_{I}\left[s(t)+n_{l}(t)\right] \cos \theta
\eeq
另一路相位相处$\frac{\pi}{2}$的信号经相同处理后得到
\beq
u_{o y}(t)=\frac{1}{2} A_{I}\left[s(t)+n_{l}(t)\right] \sin \theta
\eeq
于是被调制后信号（锁相放大器输入端信号）相对于解调信号的相位差
\beq
\theta=\arctan\frac{u_{o y}(t)}{u_{o x}(t)}
\eeq
输出信号振幅
\beq
u_{o}(t)=\sqrt{u_{\mathrm{ox}}^{2}(t)+u_{\mathrm{ov}}^{2}(t)}=\frac{1}{2} A_{I}\left[s(t)+n_{l}(t)\right]
\eeq
具有这种两互相垂直的信号输出、获得锁相放大器输入端完整信号能力的锁相放大器被称双相锁相放大器（原理如\cref{d110}），本实验用的OE1022就是双相锁相放大器。
信号的有效值
\beq
R=\sqrt{\frac{\int_{0}^{T}\left\{\left[(s(t)+n(t)) \cos \left(\omega_{0} t\right)\right]^{2}+\left[(s(t)+n(t)) \sin \left(\omega_{0} t\right)\right]^{2}\right\} d t}{T}}
\eeq
则锁相放大器输入信号的$x$、$y$分量分别为：
\bea \label{measure}
X&=R \cos \theta \approx \frac{u_{o x_{-} p p}(t)}{\sqrt{2} A_{I}} \\
Y&=R \sin \theta \approx \frac{u_{o y_{-p p}}(t)}{\sqrt{2} A_{I}}\\
R&=\sqrt{X^{2}+Y^{2}}
\eea
其中$R$、$X$、$Y$和$\theta$就是在锁相放大器屏幕上的显示值，是内部运算的结果；当噪声足够小、且$s(t)$为常量时，$R$为一常量。
\par
可见，锁相放大器用相敏检测器实现解调，并且实现既检幅又检相的功能。尽管锁相放大器已滤除了绝大部分噪声，但对同频噪声$n_l (t)$仍无能为力。并且由于放大后的信号仍包含噪声成分，等效于其相位差$\theta$也受随机噪声$n_l (t)$影响，从而存在相位噪声。在实际工作中，我们通常选择频谱中背景噪声最低的频率作为慢变信号的调制频率。\par
很多应用场合，锁相放大器的参考信号由外部提供（例如光学斩波器，信号发生器），此时叫做外部参考模式。锁相放大器的参考通道就会将外部输入的参考信号进行放大和整形，然后通过锁相环技术，在锁相放大器内部产生一个与外部输入参考信号频率相等、相位差锁定的的正弦信号。这是真的沙包，我不知道我在这里复制粘贴干嘛，可见，锁相的含义是根据外部参考，利用锁相环产生一个频率相同、相位差锁定的参考信号。这个是锁相放大器能够进行后面相敏检测运算的基础。参考通道还有另一功能，即对参考输入进行移相处理,以使各种不同相移信号的检测结果达到最佳。
\par
如果采用内部参考模式，也就是锁相放大器自带的 SINE OUT 的正弦信号作为实 验的激励源，由于内部参考信号是锁相放大器自身产生的正弦波(与 SINE OUT 信号同频不同 相)，因此频率及相位不需要锁相环获取，此时也就不用锁相;在这种情况下，由于频率及相 位是直接得到的，不需要放大和整形，测量的精度会更高。
\par

对于数字锁相放大器 Output 通道的输出信号，对应于原信号$R$或$X$、$Y$分量乘以其逻辑增益。电路实现时，LPF 的频带可以做得很窄,使锁相放大器达到较大的 SNIR。
\par
锁相放大器有几个重要参数：灵敏度，最小可测信号 MDS，时间常数，陡降，满刻度输出时的输入电平 FS，过载电平 OVL，输入总动态范围，输出动态范围，动态储备 DR，选频放大。锁相放大器的重要参数，影响着锁相放大器的工作。理解锁相放大器参数的含义及其（配置）对锁相放大器性能的影响，是正确使用锁相放大器的前提。
\begin{enumerate} 
\item 灵敏度从直观上来分析就是锁相放大器的当前显示量程，由信号通道的交流放大增益与相敏检波器解调增益决定。
\item 最小可测信号定义为输出能辨识的最小输入信号，可以理解为系统对小信号的分辨率。影响 MDS 的主
要因素有系统的内部噪声、温度漂移等，即结果会受到包括内部噪声、温度漂移的影响而产
生波动，MDS 定义为输出可以稳定在一定百分率波动下的最小输入。
\item 时间常数是低通滤波器$RC$电路的时间常数。可以简单地认为，时间常数越大，阶数越高，输出的带宽就越低，显示的测量幅度、相位等值就越稳定。然而，过大的时间常数会抹平输入信号（随时间）的变化，从而失去有用的信息。因此，在实际应用中，需要根据输入信号随时间变化的情况，协调时间常数与信噪比之间的平衡。
\item  FS 有时也称为满刻度灵敏度，它是用来表征锁相放大器测量灵敏 度的，拥有电压的量纲，与系统的总增益有关。
\item OVL 定义为锁相放大器任一级出现过载或临界过载时的输入电平。因为微弱信号检测通
常处理的是信噪比较低的输入，所以过载往往出现在噪声电压出现尖峰时。因此，可以将
OVL 理解为系统允许的最大输入噪声电压电平，即系统的最大噪声容限。应当指出，通常对
应不同的增益设置，各级放大情况将有所不同，其过载电平也会有所不同，故在指明 OVL
时，应当附带指出其 FS，这样才有意义。
\item  在给定 FS 的条件下（即给定的增益设置），锁相放大器的过载电平 OVL 与最小可测信
 号 MDS 比值的分贝数
 \beq
\text{输入总动态范围} = 20 \mathrm{lg\ } \frac{OVL}{MDS} \unit{dB}
\eeq
上文已经介绍，OVL 标称了锁相放大器的噪声容限，而 MDS 表明锁相放大器能够分辨的
最小信号。因而输入总动态范围可以理解为锁相放大器从噪声中提取有用信号的能力，即分
辨率越高，噪声容限越大，则输入总动态范围越大。
 \item 该参数定义为满刻度灵敏度 FS 与最小可测信号 MDS 比值的分贝数，即
 \beq
 \text{输出动态范围} = 20 \mathrm{lg\ } \frac{FS}{MDS} \unit{dB}
\eeq
 输出动态范围表示锁相放大器可以检测的有用输入信号的动态范围，即输入有效信号可以在该范围内波动而既不会导致锁相放大器不可分辨，也不会导致超过输出的最大范围。
 \item 动态储备 DR 定义为过载电平 OVL 与满刻度输出时的输入电平 FS 比值的分贝数
 \beq
 DR = 20 \mathrm{lg\ } \frac{OVL}{FS} \unit{dB}
\eeq
\par
锁相放大器输入信号在 PSD 处理之前需要交流放大，而在 PSD 处理之后是直流放大信
号即可。在总增益不变的情况下，如果调整交流增益增加，直流增益减小，则输入噪声经交流放大很容易使 PSD 过载，动态储备减小，同时输出的直流漂移减小。反之，如果增加直
流增益，降低交流增益，则动态储备提高，使锁相放大器具有良好的抗干扰能力，但以输出
稳定性为代价，降低了测量精度。\par
在确定的测量精度要求下，动态储备有最小值。精度要求越高，其最小值就越大。在模
拟锁相放大器中，低动态储备意味着更小的输出误差和漂移。在 OE1022 数字锁放中，高动
态储备不会增加输出误差和漂移，但是会增加输出噪声。然而，如果在 A/D 转换器前的模
拟放大器增益足够大，则其被放大的自身噪声比 A/D 转换器的噪声还大。这样，输出主要
受输入噪声影响。因此，增大模拟增益即减小动态储备并不能减小输出噪声。在分辨率要求极高的情况下，增益增大并不能提高信噪比，因此，这时可以降低增益从而提高动态储备\cite{lec}。
\item 
 选频放大对某一段频率或单一频率的信号具有突出的放大作用，而对其他频率的信号具有较强抑
 制作用的放大。
 \end{enumerate}
\emptyline
\\
\textbf{【实验前思考题】}
\begin{enumerate}
 \item {\bf 请写出拍频信号的数学表达式，并分别用时域和频域图表述出来。}
 \par
 拍频信号是将一个低频信号$\omega_l$附加在一个高频信号$\omega_h$上
 \beq
 f(t) = a_l \cos \omega_l t + a_h \cos \omega_h t
 \eeq
 傅立叶空间中的表达式为
 \beq
 f(\omega) = a_l \delta(\omega - \omega_l) + a+h \delta(\omega - \omega_h)
 \eeq
 图像如\cref{beat,fbeat}。
 \cpicn{0.6}{beat}{拍频的图像，$a_l = 2$，$\omega_l = 3$，$a_h = 5$，$\omega_h = 60$}
 \cpicn{0.6}{fbeat}{拍频在傅立叶空间中的表示}
 再例如将一个高斯波包附加到高频信号上得到的拍如\cref{gaussbeat}
 \cpicn{0.6}{gaussbeat}{高斯波包与高频信号形成拍，$\omega_l \sim 5$，$\omega_h - 60$}
 \item {\bf 噪声电压信号是个随机值，原则上，只要时间足够长，其平均值为零。如果待测信号是个恒定值，且环境噪声也是稳定的，则可以通过大量采样平均来获得待测信号值。 然而，实际测量中，信号不是常量，而是（相对）慢变量，而（环境）噪声也不稳定， 这样长时间大量采样平均的方法就不奏效，锁相放大器是个更好的解决方案。为更好的使用锁相放大器，请回答以下问题：}
 \begin{enumerate}
 \item {\bf 如何用锁相放大器检测到待测的直流信号或慢变信号？}
 \par 由于$s(t)$为慢变量，因此可以用锁相放大器按\cref{measure}测量锁相放大器的输出信号随时间的变化，输出信号$u_o(t) \propto s(t) + n_l (t)$，$n_l$为低频噪声，功率谱主要在低频部分。当然，根据输入信号的特征频率大小，进行低通滤波时的截断频率$\omega_c$也需要随着调整大小。
  \item {\bf 如用斩波器调制直流信号（如光强），被斩制后的信号仍然包含有直流分量（即平均值不为零），但该直流分量随交流信号输入锁相放大器后不会被锁相放大器检测，请从数学推导上说明。}
 \par 被斩波器调制后的直流信号仍然含有直流分量，也就是说，这个时候，被调制后的 信号可以写成如下形式
 \beq
 u_{m}(t)=x(t) \sin \left(\omega_{0} t+\theta\right)+C
 \eeq
 $C$即直流分量。之后，信号被锁相放大器放大，并经过$\sin(\omega_0 t)$解调，信号变为
 \beq
 u_{p}(t)=A_{l} x(t) \sin \left(\omega_{0} t+\theta\right) \sin \left(\omega_{0} t\right)+A_{l} C \sin \left(\omega_{0} t\right)
 \eeq
然后，信号会被低通滤波，由于$\omega_0$（参考信号的频率）相对于截止频率是高频，所以会被低通滤波器滤除掉，因此，这部分直流分量不会被锁相放大器检测到。


 \end{enumerate}
 \item {\bf 相位角以及相位差的含义是什么?锁相放大器输出的$\theta$是待测信号的相位角还是待测信号与参考信号之间的相位差?}\par
相位角和相位差是两个完全不同的概念。首先，我们来解释一下相位角。
\par
查阅百度百科\cite{phasebaidu}，了解到相位角又称相位，所以相位角其实就是我们平时经常提及的相位，它表 征了某一时刻，信号所处的状态。\par
例如，一个三角函数信号可以写为如下形式
\beq
x(t)=\operatorname{arccos}(\omega t+\varphi)
\eeq
在这个信号当中，$\omega t + \varphi$就是这个信号的相位角。当我们知道了相位角后，我们就知道了在这一时刻，信号的大小，所以我们就知道了信号所处的状态。
\par
又例如，我们可以用复数来表示一个三角函数信号，即
\beq
x^{\prime}(t)=a e^{i(\omega t+\varphi)}
\eeq
在上述表达式中，$𝑎$是复数的模，$\omega 𝑡 + \varphi$是这个复数的辐角，即我们已经提及的相位角。从这个角度，我们也可以理解，为什么我们称其为相位角。
\par
当然，很多时候，我们并不关心某个信号的相位角是多少，一方面是因为，如果信号变化的频率很快（例如光的频率是$10^{15}\unit{Hz}$量级），这个时候，整个相位角也随时间快速变化，所以我们很难准确知道信号在某一时刻的相位角是多少。另一方面，相位角总是相对的，因为初相位可以任意选取，当我们说某个信号的相位角时，其实是相对于初相位而言的。所以很多时候，测量相位角并没有很大的意义（这个是叶贤基老师光学课上讲的，此处引用）。
\par
而相位差（通常来说，比相位角更有意义），它是相对两个信号而言的，即两个信号相位角的差值被称作相位差。
\par
相位差能表征两个信号之间的关系。设$\Delta \varphi$为第一个信号与第二个信号的相位角差值：
\par
如果$\Delta \varphi>0$，那么就称第一个信号超前于第二个信号；
\par
如果$\Delta \varphi<0$，那么就称第一个信号落后于第二个信号；
\par
如果$\Delta \varphi=0$，那么就称这两个信号同相；
\par
如果$\Delta \varphi=\pm 180^{\circ}$，那么就称这两个信号反相；
\par
相位差的概念应用广泛，在各种学科中都有涉及。例如在光学\cite{ced,dtem}工程当中，我们就会用到相位差的概念。设有两束频率相同，强度也相同的光，分别表示为
\beq \label{38}
E_{1}=E_{0} e^{i\left(\omega t+\varphi_{1}\right)}
\eeq
\beq \label{39}
E_{2}=E_{0} e^{i\left(\omega t+\varphi_{2}\right)}
\eeq
现令两束光发生干涉，干涉后的光强可以写为
\beq
I=\left(E_{1}+E_{2}\right)^{*}\left(E_{1}+E_{2}\right)
\eeq
将式\cref{38}，\cref{39}代入上式，得到
\beq
I=I_{0}\left[e^{i\left(\omega t+\varphi_{1}\right)}+e^{i\left(\omega t+\varphi_{2}\right)}\right]\left[e^{-i\left(\omega t+\varphi_{1}\right)}+e^{-i\left(\omega t+\varphi_{2}\right)}\right]
\eeq
其中$I_0 = |E_0|^2$，将上式化简得到
\beq
I=I_{0}\left\{2+\left[e^{i\left(\varphi_{1}-\varphi_{2}\right)}+e^{-i\left(\varphi_{1}-\varphi_{2}\right)}\right]\right\}
\eeq
进一步化简，得到
\beq \label{43}
I=2 I_{0}+2 I_{0} \cos \Delta \varphi
\eeq
其中$\Delta \varphi=\varphi_{1}-\varphi_{2}$就是这两束光的相位差。从\cref{43}可以看见，相位差决定着干涉光强，
因而决定着两束光是相消干涉，还是相涨干涉。所以，我们可以看到，相位差是一个非常重要的概念。
\par
根据原理推导部分的式\cref{measure}，以及最初引入$\theta$的定义，我们可以很容易知道， $\theta$是待测信号与参考信号之间的相位差（具体过程在原理部分已经论述得十分详细，此处不再重复）。

 
 
 
  \item {\bf 用复数的取代三角函数，将锁相放大器的工作过程描述一遍；}
\par
待测信号$f(t) = s(t) + n(t)$，调制信号$pe^{i\omega_0 t}$，其中$p\in U(1)$是一个相位因子。调制后的信号
\beq
u_m (t) = s(t)pe^{i\omega_0 t} + n(t) pe^{i\omega_0 t}
\eeq
放大$A_I$倍得到
\beq
u_a(t) = A_Is(t)pe^{i\omega_0 t} + A_I n(t) pe^{i\omega_0 t}
\eeq
用参考信号$e^{i\omega_0 t}$解调
\bea
u_p(t) &= \mathrm{Re\ }e^{i\omega_0 t}\mathrm{Re\ }u_a(t)  + i\mathrm{Im\ }e^{i\omega_0 t}\mathrm{Im\ }u_a(t)\\
&= \frac{1}{2}p^*A_I [s(t)+n(t)] + \frac{1}{2}p A_I[s(t)+n(t)]e^{2i\omega_0t}
\eea
低通滤波得到
\beq
u_x(t) = \frac{1}{2}p^* A_I[s(t)+n_l(t)]\cos 2\omega_0 t
\eeq
用$ie^{i\omega_0 t}$调制得到的结果是
\bea
u_y(t) &= \mathrm{Re\ }e^{i\omega_0 t}\mathrm{Re\ }u_a(t)  + i\mathrm{Im\ }e^{i\omega_0 t}\mathrm{Im\ }u_a(t)\\
&= \frac{1}{2}pA_I [s(t)+n(t)] + \frac{1}{2}p^* A_I[s(t)+n(t)]\cos 2\omega_0 t
\eea

\item {\bf 用傅立叶变换描述滤波过程;帮助理解图 D1-5（时域与频域）背后的数学。【示波器带快速傅立叶变换（FFT）功能：MATH→操作（选 FFT）】}
\par
利用scipy自带的FFT函数也可以实现滤波功能，下面是将一个$f=8\unit{Hz}$的正弦信号和一个$f=20\unit{Hz}$的正弦信号加上一系列噪声混合后进行滤波的代码\cref{fft1}和最终结果\cref{fft2}。
\cpicn{0.6}{fft1}{将一个$f=8\unit{Hz}$的正弦信号和一个$f=20\unit{Hz}$的正弦信号加上一系列噪声混合后进行滤波的代码实现}
\cpicn{1}{fft2}{将一个$f=8\unit{Hz}$的正弦信号和一个$f=20\unit{Hz}$的正弦信号加上一系列噪声混合后进行滤波的最终结果}
\par
为了进一步观察滤波的作用，用\cref{fft}所示的代码作出其傅立叶变换$F(f)$如\cref{fft4}。
\cpicn{0.6}{fft}{作滤波前后频谱的代码}
\cpicn{1}{fft4}{滤波前后的频谱}
可以看到，滤波后，高频信号和高频白噪声都被大幅滤去。

\end{enumerate}

\newpage%实验记录
%\cpic{0.255}{e2}%学生信息表格
\begin{table}[H]
\centering
\begin{tabular}{|p{32mm}|p{32mm}|p{32mm}|p{32mm}|}
\hline
年级、专业： & 17级 物理学 & 组号： & 实验班6 \\ \hline
姓名： & \Stuname & 学号： & \StuID \\ \hline
日期： &  & 实验地点： & 珠海教学楼 A102 \\ \hline
学生签名： & 
\includegraphics[scale=0.09]{sign}& 评分： &  \\ \hline
日期： &  & 教师签名： &  \\ \hline
\end{tabular}
\end{table}

\begin{center}
{\LARGE{\textbf{\ExpeName}}}
\end{center}
\textbf{【实验内容、步骤、结果】}
\\
1.\textbf{锁相放大器的基本操作和基本参数——工作原理的体验}\par
1.1 \textbf{测量信号$R$、$\theta$、$X$以及$Y$值，并验证它们之间的关系}
\par
用OE1022 的内部振荡器 SINE OUT输出一个幅值$80\unit{mVrms}$，频率$1\unit{kHz}$的正弦波，并用OE1022测量，按下面步骤进行：
\begin{enumerate}
\item 断开所有与机箱连接的信号线，接入电源，打开电源开关，如果系统不是默认设置状态，按讲义18页步骤重设；
\item 在前面板上选择 REF/PHASE 子菜单，Ref.source 设置为 “Internal”;Ref.frequency 为默认值“1.000kHz”;选择进入“Sine Output”下级菜单，扫描类 型 Sweep Type 设置为“Fixed”，通过数字键盘在 Voltage 中输入“0.08”，得到幅值为 80mVrms、 频率为 1.000kHz 的正弦信号。
\item 用一条 BNC-BNC 信号线连接 OE1022 前面板 SINE OUT 输出接口和 SIGNAL IN 的 A/I
 接口，如\cref{d112}。
 \cpicn{0.6}{d112}{测量信号连线}
 \item 观察主界面中监测栏的 Overload 是否提示溢出：在显示屏最下方的Overload右边应显示NONE NONE。若前级输入溢出，则显示 Overload: INPUT NONE;若放大溢出，则显示 Overload: NONE GAIN;若同时溢出，则显示 Overload: INPUT GAIN。若发生前级溢出立即减小数字信号发生器输出幅值。按本实验的设置前级溢出应该不会发生。
 \item 调节灵敏度值。按下前面板\cref{d114} GAIN/TC 按键进入子菜单。
 \cpicn{0.6}{d114}{前面板GAIN/TC 菜单位置}
 \item 如\cref{d115}，按下软键1以选中Sensitivity功能，通过旋转旋钮调节Sensitivity 值，使测量信号值尽量满偏但不超量程，到$100\unit{mV}$左右即可。此时我们测出信号发生器发生的正弦波幅值$R=$\uline{\hspace{2cm}}，相位$\theta = $\uline{\hspace{2cm}}，参考结果见\cref{d115}。如果$\theta$非$-72.78^\circ$，可以设置REF PHASE调节。此时用示波器读数，幅值$A = $\uline{\hspace{2cm}}，对比得知$R$的含义是\uline{\hspace{6cm}}。
 \cpicn{0.6}{d115}{GAIN/TC 菜单参数显示}
 \item 主界面数据栏显示$R$、$\theta$、$X$及$Y$值。按下前面板 DISPLAY 按键进入子菜单。
\item 首先按\cref{d117}中软键 2，使其选中 Top；再按 软键 3 选中 Type，Type 区域此时高亮显示，通过调节旋钮可选择 Chart(XY 坐标)或 Bar(数字百分比)， 我们选择 Chart；再按软键 3 选 Trace，Trace 区域此时高亮显示，通过调节旋钮可选择显示$R$、$\theta$、$X$、$Y$，我们选择$\theta$。然后按下前面板右下角处的 START CONT 按键如\cref{d118}。
\cpicn{0.6}{d117}{软键选择}
\cpicn{0.6}{d118}{改变$\theta$设置}
\item 固定$R=$\uline{\hspace{2cm}}，用REF PHASE改变$\theta$值，按下前面板DISPLAY按键
进入子菜单，再按软键 1 将 Monitor 设置中从 Setting 切换为 Output，此时监测栏显示$R$、$\theta$、$X$和$Y$的值，如图 D1-20 所示。观测$X,Y$的改变，数据记录如\cref{tab:tb2}。

\begin{table}[H]
\centering

\begin{tabular}{|p{15mm}|p{15mm}|p{15mm}|p{15mm}|p{15mm}|p{15mm}|p{15mm}|p{15mm}|}
\hline
$\theta$ & $0^\circ$ & $15^\circ$ & $30^\circ$ & $45^\circ$ & $60^\circ$ & $75^\circ$ & $90^\circ$ \\ \hline
$X/{\rm mV}$ &  &  &  &  &  &  &  \\ \hline
$Y/{\rm mV}$ &  &  &  &  &  &  &  \\ \hline
\end{tabular}

\caption{$X,Y$随$\theta$的变化\label{tab:tb2}}
\end{table}
理论关系为$X= R\cos \theta$，$Y = R\sin \theta$。估算上面的测量结果与理论的吻合程度\uline{\hspace{4cm}}。


\item 改变频率值，观测$R,\theta ,X,Y$的变化，可用sweep设置参考信号，记录如\cref{tab:tb3}。

\begin{table}[H]
\centering

\begin{tabular}{|p{15mm}|p{15mm}|p{15mm}|p{15mm}|p{15mm}|p{15mm}|}
\hline
$f/{\rm kHz}$ & 1 & 3 & 10 & 30 & 100 \\ \hline
$X/{\rm mV}$ &  &  &  &  &  \\ \hline
$Y/{\rm mV}$ &  &  &  &  &  \\ \hline
$R/{\rm mV}$ &  &  &  &  &  \\ \hline
$\theta$ &  &  &  &  &  \\ \hline
\end{tabular}

\caption{$R,\theta ,X,Y$随频率的变化\label{tab:tb3}}
\end{table}

\item 设置某一参考信号频率，记录其相位差$\theta_0$，然后延长 sine out 到 signal in 的电缆，记录延长电缆后的相位差$\theta_e$，计算$\Delta \theta = \theta_e - \theta_0$如\cref{tab:tb4}。
\begin{table}[H]
\centering

\begin{tabular}{|p{15mm}|p{15mm}|p{15mm}|p{15mm}|p{15mm}|p{15mm}|}
\hline
$f/{\rm kHz}$ & 1 & 3 & 10 & 30 & 100 \\ \hline
$\theta_0$ &  &  &  &  &  \\ \hline
$\theta_e$ &  &  &  &  &  \\ \hline
$\Delta \theta$ &  &  &  &  &  \\ \hline
\end{tabular}

\caption{观察一米电缆对相位差的影响\label{tab:tb4}}
\end{table}
操作时{\red 不要}先将$\theta_0$全部记录完，再记录$\theta_e$，必须在同一频率下记录$\theta_0$和$\theta_e$后再记录下一个频率的数据。



\end{enumerate}
\emptyline
\par
1.2 \textbf{不同时间常数、陡降、动态储备下观测滤波器效果}
\par
现在我们举例使用 OE1022 的内部振荡器 SINE OUT 产生一个幅值为$50\unit{mVrms}$、频率为$1\unit{kHz} $的正弦波，用 OE1022 进行测量并在 Fast channel out 进行输出，本操作指南将简单演示如何调节时间常数（TC）、滤波器陡降（Slope）及动态储备（Reserve），通过锁相放大器的 Fast out 功能及示波器观测滤波器效果。步骤如下:
\begin{enumerate}
\item 断开所有与机箱连接的信号线，接入电源，打开电源开关，此时系统处于默认设置状态。
在前面板上选择 REF/PHASE 子菜单，Ref.source 设置为“Internal”;Ref.frequency 为默认 值“1.000kHz”;选择进入“Sine Output”下级菜单，扫描类型 Sweep Type 设置为“Fixed”，通过数字键盘在 Voltage 中输入“0.05”，得到幅值为 50mVrms、频率为 1.000kHz 的正弦信号。
\item 用 BNC-BNC 信号线连接 OE1022 前面板 SINE OUT 输出接口和 SIGNAL IN 的 A/I 接口， OE1022 前面板 CH1 接口和示波器输入接口。
\item 调节灵敏度值，本例中数字信号发生器输出有效值为$50\unit{mVrms}$的正弦波，相位（phase）设置为$60^\circ$，灵敏度设置为$100\unit{mV}$即可。
\item 按下 OE1022 前面板的 GAIN/TC 按键，调节时间常数和陡降，设置时间常数为$30\unit{\mu s}$，陡 降为$6\unit{dB/oct}$。
\item 按下 OE1022 前面板 CHANNEL OUTPUT 按键，进入子菜单，在 Output 选项中选择 CH1， 在 Speed 选项中选择 Fast，Source 选项默认为 R(Fast 模式下输出源可选择 R、X、Y， Slow 模式下输出源可选择 R、X、Y、Rh1、Xh1、Yh1、Rh2、Xh2、Yh2、Noise、$\theta$、$\theta$h1、 $\theta$h2)，Offset\&Expand 中默认 Offset 为 0，Expand 为 1，如图\cref{d121}所示。
\cpicn{0.6}{d121}{Channel out 子菜单设置图}
\item 在示波器上调节时基、输入通道设置等，可观测到在$30\unit{\mu s}$时间常数和$6\unit{dB/oct}$陡降下
channel out 输出$R$值波形，如图\cref{d122}所示：
\cpicn{0.6}{d122}{ Fast 模式下 CH1 输出$R$信号波形（$30\unit{\mu s},\  6\unit{dB/oct}$）}
示波器所测量的 CH1 输出信号频率与设定信号频率是什么关系？时间常数大小对此有何影响？为什么？\footnote{参考实验讲义24页下方。}
\vspace{5cm}

\item 在 OE1022 的 Channel out 菜单中设置 Source 为 X、Y，在示波器上观测。
\item 在 OE1022 的 Gain/TC 菜单中，对时间常数进行改变，时间常数可选择 $10\unit{\mu s}$、$30\unit{\mu s}$、$100\unit{\mu s}$、 $300\unit{\mu s}$......$10\unit s$、$30\unit s$，通过增大时间常数的值能够使系统的输出端更稳定，也能减轻输入端 噪声对输出端的影响。时间常数除了对系统的稳定性和精度有影响外，还会影响系统的响 应时间，时间常数还决定噪声测量时的等效噪声带宽(ENBW)。在此特别说明一下，等效 噪声带宽指的并不是滤波器的-3 dB 带宽，它指的是对高斯噪声10的有效带宽。观察示波器波形，时间常数为\uline{\hspace{2cm}}个信号周期时，才能观察到稳定的直流输出。
定性记录并简单分析波形变化：\footnote{参考讲义26页。}
\vspace{5cm}
\item 在 OE1022 的 Gain/TC 菜单中，对滤波器陡降进行改变。在同样的测量准确度下，分别选择 6dB/oct、12dB/oct、18dB/oct、24dB/oct 四档，在 TC 不大于 30$\unit{\mu s}$的情况下，观察并定性记录示波器波形的变化：
\vspace{5cm}


\end{enumerate}
\emptyline
\par
1.3 \textbf{制定方案，观测不同时间常数下，输出信号对输入信号变化的响应}
\par
改变 SINE OUT 幅值设定，用“START CONT”激活 chart 图以观察改变前后的曲线来确定其输出响应，并记录图形。
\par
在 OE1022 的 Gain/TC 菜单中，对滤波器陡降进行改变。在同样的测量准确度下，分别选择 6dB/oct、12dB/oct、18dB/oct、24dB/oct 四档，观察在输入信号改变时，输出信号对输入信号变化的响应；并理解滤波器陡降的作用。图形由示波器屏幕图像给出，实验过程中定性观测到的现象记录如下：
\vspace{5cm}
\emptyline
\par
1.4 \textbf{外部输入参考信号}
\par
用信号发生器产生正弦信号作为外部参考信号（外部参数信号要求峰值 100mV 以上），并用三通将信号用 BNC 线分别接入 REF IN，和 signal in 的 A/I，在锁相放大器参数设置 Ref.source->external。
\par
外部参考信号峰值$V_p = $\uline{\hspace{2cm}}，重复1.1部分的测量，记录数据如\cref{tab:tb6}。

\begin{table}[H]
\centering

\begin{tabular}{|p{15mm}|p{15mm}|p{15mm}|p{15mm}|p{15mm}|p{15mm}|p{15mm}|p{15mm}|}
\hline
$\theta$ & $0^\circ$ & $15^\circ$ & $30^\circ$ & $45^\circ$ & $60^\circ$ & $75^\circ$ & $90^\circ$ \\ \hline
$X/{\rm mV}$ &  &  &  &  &  &  &  \\ \hline
$Y/{\rm mV}$ &  &  &  &  &  &  &  \\ \hline
\end{tabular}

\caption{$X,Y$随$\theta$的变化\label{tab:tb6}}
\end{table}
理论关系为$X= R\cos \theta$，$Y = R\sin \theta$。估算上面的测量结果与理论的吻合程度\uline{\hspace{4cm}}。改变频率值，观测$R,\theta ,X,Y$的变化，可用sweep设置参考信号，记录如\cref{tab:tb7}。

\begin{table}[H]
\centering

\begin{tabular}{|p{15mm}|p{15mm}|p{15mm}|p{15mm}|p{15mm}|p{15mm}|}
\hline
$f/{\rm kHz}$ & 1 & 3 & 10 & 30 & 100 \\ \hline
$X/{\rm mV}$ &  &  &  &  &  \\ \hline
$Y/{\rm mV}$ &  &  &  &  &  \\ \hline
$R/{\rm mV}$ &  &  &  &  &  \\ \hline
$\theta$ &  &  &  &  &  \\ \hline
\end{tabular}

\caption{$R,\theta ,X,Y$随频率的变化\label{tab:tb7}}
\end{table}
根据测量结果，回答：对于 OE1022 锁相放大器的信号处理质量，用内部参考信号好还是外部参考信号好?
\vspace{2cm}
\emptyline
\\
2.\textbf{强噪声背景下的弱信号检测——应用体验}
\par
本部分实验的主要目的是理解锁相放大技术的基本原理及使用技能，掌握利用相关性原理提取深埋于强噪声中的微弱信号，原理图如\cref{d127}。
\cpicn{0.6}{d127}{强噪声背景检测弱信号实验原理图}
教学实验箱自带 100m Vrms 的白噪声发生器，实验箱内部电路利用双极性晶体管散粒噪声 的功率特性来产生白噪声；通过拨码器选择可以得到 100 mVrms、10 mVrms、1 mVrms 的白噪声信号，分别对应 0dB（1 倍）、20dB（10 倍）、40dB（100 倍）的衰减。由锁相放大器 OE1022 的 SINE OUT 以及外部的 80dB（10000 倍）信号衰减器，可以得到 100nVrms 到 5Vrms 的正弦信号。通过运放把两路信号相加得到不同信噪比的信号，然后由锁相放大器对此信号进行提取测量。
\par
实验按如下步骤进行：
\begin{enumerate}
\item 用示波器观察教学实验箱噪声发生器输出的噪声是否白噪声？记录时域和频域的图谱来说明。
\item 不同的噪声水平下，考察 OE1022 如何选用合适的动态储备测量信号?
\par
在 OE1022 的 Gain/TC 菜单中，对动态储备进行设置，可设置为 Low、High、Normal 三档之一。动态储备表示锁相放大器对噪声容忍程度的大小，OE1022 动态储备可大于 100 dB，高的动态储备会产生输出噪声和漂移，当动态储备较高时，由于模数转换器的噪声存在导致输出误差增加。由于所有的信号源都存在本底噪声，固在 PSD 提取信号过程中就会掺杂着噪声，如果噪声很大，在高动态储备测量中就会产生较大的输出误差。如果 外部噪声较小，则其输出主要是受 OE1022 自身噪声影响。这时可以通过降低动态储备和直流增益可以减小输出误差。因此，在实际应用中应尽量使用低动态储备。默认配置下系统会根据输入信号的 R 值自动调整所需要的最小动态范围。观察示波器波形变化，给出结论如下：
\vspace{5cm}
\item 比较 OE1022 与示波器对不同信号强度和信噪比的测量结果：
\begin{enumerate}
\item 使用 OE1022 产生频率 1 kHz，幅值为 100 mVrms（0.282Vpp）的正弦波信号;
\item 对锁相放大器 OE1022 进行以下设置：
	\begin{enumerate}
	\item 进入 INPUT FLTERS 菜单，设置 Source 为 A;
	\item 进入 GAIN TC 菜单，设置 Sensitivity 为 500 mV，Reserve 为 Normal，Time Constant 为 1 s，Filter dB/oct 为 24 dB；
	\item 进入 REF PHASE 菜单，设置 Ref. source 为 Internal；
	\end{enumerate}
\item 使用 BNC-BNC 信号线连接 OE1022 的“SINE OUT”接口与实验箱本实验框图中的
“VIN”接口；
\item 通过分线器，使用 BNC-BNC 信号线将实验箱本实验框图中的“VOUT”接口分别连接
到 OE1022 的“A/I”接口、和示波器输入端；
\item 读取 OE1022 测到的 R 值，即为被噪声信号淹没的正弦信号有效值；
\item 改变 OE1022 产生正弦波有效值，在不同信噪比下重复上述测量；将实验测得波形图、R 值记录在表\cref{tab:tb21}中，作实验结果分析。
\end{enumerate}
\end{enumerate}
注意：用三通将输出信号分别连接到锁相放大器的输入端(A/I)和示波器的输入端；在用示波器观察、记录时，请注意屏幕水平轴范围应达到待观察信号的若干周期。
\begin{table}[H]

\cpic{0.42}{tb21}

\caption{强噪声背景检测弱信号实验结果\label{tab:tb21}}
\end{table}
\newpage
\textbf{【现场预习思考题】}
\begin{enumerate}
\item {\bf 噪声有哪些类型？一般测量对象本身的噪声是哪来的?它有什么特征？}\par
 噪声有三类：热噪声、散粒噪声和$1/f$噪声\cite{mpl}；这些噪声都是由元器件内部电子运动的涨落现象引起的。\par
噪声的特征是无规、随机的变化。其每个傅立叶模$f(\omega)$是一个均值为0，方差为$P(\omega)$（功率谱）的随机变量。对热噪声和散粒噪声，它们都属于白噪声\cite{white_noise}，功率谱
\beq
P(\omega) = p_0
\eeq
为常数；$1/f$噪声的功率谱与频率成反比
\beq
P(\omega) = \frac{P}{\omega}
\eeq
在低频处的影响非常显著。
\item {\bf 描述用教学实验箱配制不同信噪比信号的原理;实验箱生成的噪声与实际噪声是否一样？（请通过实验测量后再回答是否为白噪声）}\par
教学试验箱通过拨码器选择可以得到不同信噪比的白噪声信号。\par
实验箱利用双极性晶体管散粒噪声的功率特性来产生白噪声，这种白噪声是由散粒噪声、热噪声、配分噪声和$1/f$噪声组成的。晶体管基区和各项电阻上载流子的不规则热运动产生的电流涨落即为热噪声。晶体管中少数载流子通过发射极-基极结注入到基区时，少数载流子的数目和速度都有涨落，引起通过结的电流的微小变化。同时，少数载流子在基区内的不规则运动也将引起电流涨落，这些都产生晶体管的散粒噪声。在晶体管基区中，发射极电流的一部分变为集电极电流，另一部分变为基极电流，有一个由空穴-电子复合作用而定的电流分配系数。复合现象本身同样受到热起伏效应的影响，因此分配系数不是恒定的。它的微小变化也会引起集电极电流的起伏，这就是晶体管的配分噪声。在晶体管噪声频谱中，低频时噪声急剧上升,呈$1/f$关系，即其功率谱
\beq
P(\omega) \sim \frac{1}{\omega}
\eeq
随工艺条件、表面处理和环境气氛等的不同，取$f = 1-2$之间,故常称为$1/f$噪声。
\par
实验中用示波器直接测量实验箱的输出如，可以看到，其生成的噪声与实际中遇到的噪声无显著差异。
\item {\bf 随着信噪比的降低，信号$R$、相位差$\theta$的测量值会有什么影响？（用实验结果回答）}
\par
随着信噪比降低，信号$R$随之降低，相位差$\theta$减小。
\end{enumerate}






\newpage

\textbf{【实验过程中遇到问题记录】}\par
%\cpic{0.81}{none}
%生成最终报告时将上面内容全部删除或注释（用\iffalse \fi），将扫描得到的实验报告保存为Record.pdf在LaTeX model for GPL下，将下行命令的注释号删去。注意根据实际页数调整pages参数。
%\includepdf[pages=1-3]{Record}


\newpage%分析与讨论
%\cpic{0.255}{e3}%学生信息表格
\begin{table}[H]
\centering
\begin{tabular}{|p{32mm}|p{32mm}|p{32mm}|p{32mm}|}
\hline
年级、专业： & 17级 物理学 & 组号： & 实验班6 \\ \hline
姓名： & 高寒 & 学号： & 17353019 \\ \hline
日期： & \today &  &  \\ \hline
评分： &  & 教师签名： &  \\ \hline
\end{tabular}
\end{table}
\begin{center}
{\LARGE\textbf{{\ExpeName}}}
\end{center}
\textbf{【分析与讨论】}\\

\newpage

\bibliographystyle{IEEEtran}
\bibliography{model.bib}

\end{document}